\documentclass[11pt,a4paper]{article}
\usepackage[utf8]{inputenc}
\usepackage[T1]{fontenc}
\usepackage{geometry}
\usepackage{amsmath,amssymb,amsthm}
\usepackage{booktabs}
\usepackage{array}
\usepackage{xcolor}
\usepackage{listings}
\usepackage[hidelinks]{hyperref}
\usepackage{mdframed}
\usepackage{enumitem}
\usepackage{microtype}
\usepackage{graphicx}

\geometry{margin=1.1in}

\definecolor{rustbrown}{HTML}{704214}
\definecolor{tsblu}{HTML}{1A4A8A}
\definecolor{schemagreen}{HTML}{1A6B3A}
\definecolor{codegray}{HTML}{F5F5F5}
\definecolor{ruleblue}{HTML}{2C5F8A}

\lstset{
  basicstyle=\ttfamily\small,
  backgroundcolor=\color{codegray},
  frame=single,
  framerule=0pt,
  breaklines=true,
  columns=flexible,
  keepspaces=true,
  xleftmargin=8pt,
  xrightmargin=8pt,
}

\newmdenv[
  backgroundcolor=ruleblue!8,
  linecolor=ruleblue,
  linewidth=1.5pt,
  leftmargin=0,
  rightmargin=0,
  innerleftmargin=10pt,
  innerrightmargin=10pt,
  innertopmargin=8pt,
  innerbottommargin=8pt,
]{rulebox}

\theoremstyle{definition}
\newtheorem{theorem}{Theorem}
\newtheorem{definition}{Definition}
\newtheorem{corollary}[theorem]{Corollary}

\title{\textbf{The Type System as Press Release}\\[6pt]
\large $A = \mu(\mathcal{O}^*)$ and the Instruments of Ontological Closure}
\author{Sean Chatman\\
\small ChatmanGPT / pictl project}
\date{April 2026}

\begin{document}
\maketitle

\begin{abstract}
The \textit{Chatman Equation} $A = \mu(\mathcal{O}^*)$ asserts that every software artifact $A$ is a lawful transformation $\mu$ of an operationally complete source ontology $\mathcal{O}^*$. This demands that code, schemas, documentation, tests, and generated interfaces be deterministic projections from one governing semantic source---not improvisations from partial observation. This paper surveys the practical instruments by which $\mathcal{O}^*$ is encoded, propagated, and enforced across compiler, runtime, human reader, and cross-language boundary: Rust's type system with \texttt{rustdoc} and doctests; TypeScript's static types paired with Zod for runtime enforcement; Protocol Buffers with Buf for cross-language contracts; JSON Schema and OpenAPI 3.1; and the Serde/Schemars bridge for Rust-to-schema emission. We introduce the \textit{Blue River Dam} principle---upstream semantic control over all admissible downstream engineering work---and argue that the strongest tool stacks are those that preserve $\mathcal{O}^*$ across the most system boundaries with the least semantic drift. A five-dimension maturity matrix ranks each instrument. We conclude with a 5~Whys and Root Cause Analysis applied to the \texttt{pictl}$\to$\texttt{bcinr} transition, demonstrating that the Working Backwards architecture's true power is not avoiding failure but converting exploration results into rational press release revision---and that this revision is itself provable by mathematical oracle.
\end{abstract}

\tableofcontents
\newpage

%% ────────────────────────────────────────────────────────────────
\section{The Chatman Equation: $A = \mu(\mathcal{O}^*)$}

In a naïve software development model, artifacts emerge bottom-up from local decisions: a function is written because a test needed something to call; a schema is written because a database column needed a name. Correctness is a post-hoc property, verified after construction.

The Chatman Equation inverts this. It asserts:

\begin{definition}[Chatman Equation]
Every correct artifact $A$ is the unique image of an operationally complete source ontology $\mathcal{O}^*$ under a lawful transformation function $\mu$:
$$A = \mu(\mathcal{O}^*)$$
\end{definition}

The ontology $\mathcal{O}^*$ is not a database schema or a type annotation in isolation. It is the \emph{complete semantic specification} of the intended system: the union of all domain commitments that, when fully resolved, determine every artifact without ambiguity. $\mu$ is the transformation---a code generator, a compiler, a type checker, a documentation renderer---that projects $\mathcal{O}^*$ into a particular artifact form.

The relevance to tool selection follows directly. The question is not ``which tool is most popular?'' It is:

\begin{rulebox}
\textbf{Which tool most faithfully preserves $\mathcal{O}^*$ across the most system boundaries with the least semantic drift?}
\end{rulebox}

A type annotation that disappears at runtime (TypeScript without runtime validation) only partially preserves $\mathcal{O}^*$. A Protobuf IDL that is used to generate code in five languages preserves it very broadly. A docstring that describes behavior not verified by any test preserves almost none of it.

This is the governing lens for everything that follows.

%% ────────────────────────────────────────────────────────────────
\section{Before Ontology Closure: Probabilistic vs.\ Deterministic Outputs}

A critical threshold in the Chatman Equation is \emph{ontology closure}. Before closure, some commitments in $\mathcal{O}^*$ are ambiguous or absent; $\mu$ must fill gaps through interpretation, convention, or default. Outputs are probabilistic in the sense that different developers, generators, or tools will produce different artifacts from the same incomplete spec.

After closure---when every domain commitment is explicitly stated, typed, and constrained---$\mu$ becomes deterministic. The ontology \emph{schedules} artifact generation rather than guiding it.

\begin{theorem}[Closure Determines Delivery]
Let $\Omega$ be an oracle that evaluates whether a realized system $R$ satisfies the press release $P$. Then:
$$\text{Schedulable}(P) \propto \text{Strength}(\Omega)$$
The stronger the oracle (Rank~1 mathematical theorem $>$ Rank~4 statistical trend), the more precisely $P$ can be converted into a delivery schedule rather than a wish.
\end{theorem}

This is why the \texttt{bcinr} library was chosen as the post-\texttt{pictl} target. Its domain---branchless integer algorithms---admits Rank~1 oracles:
$$\forall a, b \in \mathbb{Z}_2^{32},\quad \texttt{select\_u32}(\texttt{0xFFFF\_FFFF}, a, b) = a$$
Such a statement is not a test; it is a theorem. The type system plus a doctest that executes this assertion turns the theorem into a continuously enforced delivery constraint.

%% ────────────────────────────────────────────────────────────────
\section{Blue River Dam: Upstream Control Over Engineering Flow}

Blue Ocean Strategy asks: how do we find uncontested market space? The software equivalent is analogous but insufficient for frontier systems. Knowing where the ocean is does not tell you what governs the flow of engineering work toward it.

We introduce a stronger framing.

\begin{definition}[Blue River Dam]
A \emph{Blue River Dam} is the upstream governing artifact that determines what engineering work is admissible downstream. In frontier software development, the dam is the press release / plan file / ontology layer. Whoever controls that upstream definition controls what work exists at all.
\end{definition}

The distinction from Blue Ocean:
\begin{itemize}
\item \textbf{Blue Ocean:} Find empty water (uncontested problem space).
\item \textbf{Blue River Dam:} Control the headwaters (govern what work is even admitted).
\end{itemize}

In practical terms, the Blue River Dam is realized by the combination of instruments in this paper:
\begin{enumerate}
\item The \emph{type system} controls what values are admissible at the language level.
\item The \emph{schema / IDL} controls what data is admissible across process and language boundaries.
\item The \emph{doctest} controls what examples are admissible in documentation.
\item The \emph{ontology / plan file} controls what features are admissible in the product.
\end{enumerate}

If any of these layers is absent or inconsistent, the river---the stream of tickets, APIs, implementations---detaches from the vision and flows toward wherever it encounters least resistance.

%% ────────────────────────────────────────────────────────────────
\section{The Instrument Taxonomy}

We organize the practical tools into five categories, ordered by abstraction layer from most fundamental (language semantics) to most portable (cross-system contracts).

\subsection{Language-Level Semantics}

\paragraph{Rust type system.}
Rust provides the strongest single-language type system in mainstream use. Traits encode behavior contracts; lifetimes encode ownership and aliasing invariants; \texttt{Result<T,E>} makes failure paths explicit. The type checker is a continuous, zero-cost Rank~1 oracle at the language boundary.

\paragraph{TypeScript type system.}
TypeScript provides rich structural typing with generics, mapped types, and discriminated unions. \textbf{Critical limitation:} TypeScript type annotations are erased from emitted JavaScript. The type system governs author-time reasoning; it does not govern runtime data.

\paragraph{JSDoc-annotated JavaScript via TypeScript.}
TypeScript can check JavaScript files annotated with JSDoc type comments. This is a transitional or migration-path option. It preserves less precision than native TypeScript and cannot represent advanced type constructs.

\subsection{Documentation Extraction and Drift Control}

\paragraph{Rustdoc + doctests.}
\texttt{rustdoc} generates documentation from source comments via \texttt{///} doc-comment syntax. The critical property: examples in doc comments are compiled and executed as tests by \texttt{cargo test}. This closes the drift loop between documentation and implementation at the source level. There is no separate documentation maintenance burden; documentation that diverges from implementation fails the test suite.

\paragraph{TSDoc + TypeDoc.}
TSDoc standardizes documentation comment syntax for TypeScript tools. TypeDoc generates HTML and JSON documentation from TypeScript exports. Documentation examples are \emph{not} executed as tests by default, which means they can silently diverge from implementation.

\subsection{Runtime Validation / Schema Declaration}

\paragraph{Zod.}
Zod is TypeScript-first runtime validation with schema-inferred static types. The same schema object produces both the TypeScript type annotation and the runtime parse/validate logic. This closes the TypeScript compile-runtime split.

\paragraph{JSON Schema.}
JSON Schema is a language-neutral vocabulary for describing and validating JSON data. It is widely supported across tools, frameworks, and languages. Precision is lower than a full language type system but portability is high.

\subsection{API / IDL Contract Publication}

\paragraph{OpenAPI 3.1.}
OpenAPI defines a standard interface for HTTP APIs readable by both humans and tools. OpenAPI 3.1 aligns with JSON Schema 2020-12. It enables documentation generation, client SDK generation, gateway configuration, and design review from one spec file.

\paragraph{Protocol Buffers + Buf.}
Protobuf is a language-neutral, platform-neutral serialization format and IDL with code generation across languages. Buf adds linting, breaking-change detection, and schema registry capabilities to Protobuf workflows. This combination is the most operationally mature choice when contracts must survive across team, repository, and language boundaries.

\subsection{Code Generation Bridges}

\paragraph{Serde + Schemars (Rust).}
Serde provides efficient generic serialization/deserialization via derive macros. Schemars derives JSON Schema from Rust types, reflecting the JSON representation those types produce under \texttt{serde\_json}. Together they project Rust domain types into exchange-portable schemas without maintaining a separate schema source.

\paragraph{tonic (Rust gRPC).}
\texttt{tonic} is a high-performance Rust gRPC implementation focused on interoperability. It realizes Protobuf contracts at the Rust transport layer. It is below Protobuf/Buf in the ontological stack; it realizes the contract rather than defining it.

%% ────────────────────────────────────────────────────────────────
\section{Maturity Matrix}

We score each instrument on five dimensions, each rated 1--5:

\begin{description}[leftmargin=2cm, labelwidth=2cm]
  \item[Precision] How much domain semantics can be expressed and enforced.
  \item[Enforcement] Whether violations are caught at compile time, validation time, or test time.
  \item[Interop] Whether the contract survives outside one language or runtime.
  \item[Drift Resistance] Whether documentation, examples, and contracts remain synchronized with implementation.
  \item[Maturity] Ecosystem depth, long-term stability, and operational track record.
\end{description}

\medskip
\begin{center}
\begin{tabular}{@{}l c c c c c c@{}}
\toprule
\textbf{Instrument} & \textbf{Prec} & \textbf{Enf} & \textbf{Interop} & \textbf{Drift} & \textbf{Mat} & \textbf{Mean} \\
\midrule
Rust types + rustdoc + doctests & 5 & 4 & 2 & 5 & 5 & \textbf{4.6} \\
Protobuf + Buf                  & 4 & 4 & 5 & 5 & 5 & \textbf{4.6} \\
TypeScript + Zod                & 4 & 5 & 3 & 4 & 5 & \textbf{4.2} \\
JSON Schema                     & 3 & 5 & 5 & 4 & 5 & \textbf{4.4} \\
OpenAPI 3.1                     & 3 & 4 & 5 & 4 & 5 & \textbf{4.2} \\
Serde + Schemars                & 4 & 4 & 4 & 4 & 4 & \textbf{4.0} \\
tonic                           & 3 & 4 & 4 & 3 & 4 & \textbf{3.6} \\
TSDoc + TypeDoc                 & 3 & 1 & 2 & 3 & 4 & \textbf{2.6} \\
JSDoc + TS checking             & 2 & 1 & 2 & 2 & 5 & \textbf{2.4} \\
Docstrings alone                & 1 & 1 & 1 & 1 & 5 & \textbf{1.8} \\
\bottomrule
\end{tabular}
\end{center}

\medskip

\noindent
\textbf{Reading the matrix.}
Rust types + rustdoc + doctests and Protobuf + Buf share the highest mean score for opposite reasons: Rust scores highest on local precision and drift resistance; Protobuf scores highest on interoperability and boundary-crossing contract durability. TypeScript without Zod would score 3.0 (Enforcement drops to 2 because runtime ingress is ungoverned). Docstrings alone score 1.8---they have ecosystem ubiquity but enforce nothing.

%% ────────────────────────────────────────────────────────────────
\section{Stack Rankings by Use Case}

\subsection{Rust-first Core Library}
\begin{enumerate}
\item Rust type system
\item \texttt{rustdoc} + doctests
\item Serde derive macros
\item Schemars (JSON Schema emission at boundaries)
\item Protobuf / OpenAPI at service boundaries if needed
\end{enumerate}

\noindent $\mathcal{O}^* \to \text{Rust types} \to \text{docs/tests} \to \text{exchange formats}$. This minimizes duplicate declarations; the source type drives all downstream forms.

\subsection{TypeScript Application or SDK}
\begin{enumerate}
\item TypeScript types
\item Zod (runtime ingress validation)
\item TSDoc annotations
\item TypeDoc for API surface documentation
\item OpenAPI 3.1 / JSON Schema for published interfaces
\end{enumerate}

\noindent $\mathcal{O}^* \to \text{TS static model} \to \text{Zod runtime parser} \to \text{published contract}$. Without Zod, you only control author-time reasoning, not the ingress truth.

\subsection{Polyglot Microservices}
\begin{enumerate}
\item Protobuf + Buf (canonical semantic source)
\item Language-native generated types from proto
\item Service runtime (tonic, gRPC-java, etc.)
\item Human documentation from proto comments
\end{enumerate}

\noindent $\mathcal{O}^* \to \text{IDL} \to \text{generated clients/servers} \to \text{runtime}$.

\subsection{HTTP APIs for Broad Consumption}
\begin{enumerate}
\item Domain model (TS types or Rust types)
\item JSON Schema-aligned schema layer
\item OpenAPI 3.1 (publication artifact)
\item Generated SDKs and docs
\item Runtime validator (Zod, Pydantic, etc.)
\end{enumerate}

\noindent OpenAPI is the publication surface, not the deepest semantic source. The domain model governs it.

\subsection{The Universal Anti-Pattern}
\begin{lstlisting}
// WRONG: docstring is the canonical source
/// Returns the sum of a and b
fn add(a: u32, b: u32) -> u32 {
    a - b  // compiles fine, docstring says nothing enforced
}
\end{lstlisting}

\noindent The docstring cannot catch the bug. The type system alone cannot catch the semantic error. Only a doctest that asserts \texttt{add(2, 3) == 5} closes the loop.

%% ────────────────────────────────────────────────────────────────
\section{5~Whys and Root Cause Analysis: pictl $\to$ bcinr}

\subsection{Why This RCA Proves the Exploit Theory}

In the companion paper \textit{``Working Backwards in Code''} we introduced the Explore/Plan/Write agent architecture as a computational Working Backwards engine. The pictl adversarial audit was the Explore phase: 41 algorithms claimed, 0 production-ready. This section applies 5~Whys to that result and shows how it produces the bcinr press release---not as consolation but as rational exploitation of the information gained.

The \emph{Exploit theory} states: once Exploration has measured the gap precisely enough to identify structural root cause, the correct response is not incremental repair but press release revision toward a target whose gap is closable with mathematical oracles. The bcinr transition is the proof.

\subsection{5~Whys Chain}

\paragraph{Why 1: Why were 0 of 41 pictl algorithms production-ready?}
Because 24 crashed on invocation with \texttt{RefCell} borrow panics, 2 returned wrong types, and the remainder fell into stub or aspirational categories. None met the Tier~0 standard (returns correct typed output, matches declared performance, passes adversarial inputs).

\paragraph{Why 2: Why did the RefCell borrow panics occur?}
Because the architecture used \texttt{RefCell<Mutex<HashMap<...>>>} as the shared state primitive. \texttt{RefCell} enforces single-owner semantics at runtime; \texttt{Mutex} is a thread-safety primitive. Their combination is architecturally incoherent in single-threaded WASM, where the thread-safety guarantee is unnecessary but the runtime borrow-checking overhead is real and causes panics when borrows are not manually released between calls.

\paragraph{Why 3: Why was this architectural defect not detected earlier?}
Because the development process used sequential ticket exploration: one algorithm at a time, pass/fail per ticket. Sequential exploration surfaces \emph{local} facts. It cannot distinguish between ``this algorithm has a fixable bug'' and ``every algorithm shares the same architectural flaw.'' The architectural pattern---that all 24 crashes were caused by the same root type---only became legible when all 41 algorithms were audited simultaneously.

\paragraph{Why 4: Why was sequential ticket exploration chosen over parallel architectural audit?}
Because ticket systems optimize velocity (tasks/sprint) rather than gap closure ($-d\Delta/dt$). Each ticket appeared independently completable. There was no mechanism to aggregate crash signals into an architectural inference. The press release (41 algorithms registered) measured gap completion in units of registered names, not units of semantic correctness.

\paragraph{Why 5: Why did the press release measure in registered names rather than semantic correctness?}
Because no oracle was specified when the registry was created. A name is the weakest possible oracle: the presence of a string in a registry. Without stronger oracles---typed return values, performance bounds, adversarial test pass rates---the gap metric was underpowered. Any string could satisfy it.

\subsection{Root Cause}

\begin{rulebox}
\textbf{Root Cause:} The pictl press release committed to algorithm names without committing to mathematical oracles. This made the gap metric ($\Delta$) insensitive to the difference between ``exists'' and ``works.'' Sequential exploration could not aggregate crash signals into the architectural inference. The result was an architecture (RefCell/Mutex in WASM) that was fundamentally incompatible with the declared target, invisible until audited in full.
\end{rulebox}

\subsection{The Exploit: bcinr as Rational Press Release Revision}

The 5~Whys chain produces the corrective specification directly:

\begin{center}
\begin{tabular}{@{}l l@{}}
\toprule
\textbf{pictl defect} & \textbf{bcinr corrective commitment} \\
\midrule
Weak oracle (name existence) & Rank-1 mathematical theorems as oracle \\
Architectural WASM incompatibility & \texttt{no\_std}, no \texttt{Mutex}, no \texttt{RefCell} \\
Sequential exploration & Ontology-driven parallel generation (unRDF) \\
Registry without types & RDF ontology with typed API symbols \\
Docstrings not enforced & rustdoc doctests executed by \texttt{cargo test} \\
Gap metric: names & Gap metric: doctest pass rate \\
\bottomrule
\end{tabular}
\end{center}

Every bcinr architectural decision is a direct consequence of the pictl 5~Whys chain. This is what distinguishes press release revision from retreat: the new press release is narrower precisely because the exploration revealed which structural properties made the old one undeliverable.

The \textit{Exploit} is not ``try again.'' It is ``commit to a target whose gap can be closed with instruments of higher oracle rank.''

%% ────────────────────────────────────────────────────────────────
\section{The Canonical Projection Order}

Synthesizing the maturity matrix, the Blue River Dam principle, and the 5~Whys analysis, we arrive at the canonical projection order for press release instruments:

$$\mathcal{O}^* \;\xrightarrow{\mu_1}\; \text{Types} \;\xrightarrow{\mu_2}\; \text{Doctests} \;\xrightarrow{\mu_3}\; \text{Schemas} \;\xrightarrow{\mu_4}\; \text{Published Contract}$$

Each arrow is a projection from the source. The most important architectural rule is:

\begin{rulebox}
\textbf{Do not let docstrings be the canonical source of truth.}\\[4pt]
Prefer: \quad $\mathcal{O}^*\;{\to}\;\text{generated docs}\;{\to}\;\text{validated examples}\;{\to}\;\text{published contract}$
\end{rulebox}

In bcinr, this order is realized as:
\begin{enumerate}
\item $\mathcal{O}^*$ = \texttt{ontology/bcinr.ttl} (RDF Turtle, 12 algorithm families, typed symbols)
\item \texttt{api/*.rs} = generated public surface (via \texttt{unrdf sync}, overwritten on each generation)
\item \texttt{logic/*.rs} = handwritten hot loops (bootstrap-once, human-owned implementation)
\item Doctests in \texttt{logic/*.rs} = executable oracle assertions, executed by \texttt{cargo test}
\item JSON Schema via Schemars = export surface for downstream consumers
\end{enumerate}

This is the Chatman Equation instantiated in a complete project:
$$A_{\text{bcinr}} = \mu(\texttt{bcinr.ttl})$$

%% ────────────────────────────────────────────────────────────────
\section{Implications for Agent-Driven Development}

\subsection{Plan File = Press Release = $\mathcal{O}^*$}

In the Explore/Plan/Write agent architecture, the plan file is the computational press release. It must contain:
\begin{itemize}
\item The desired state $P$ (what the system is claimed to do)
\item The measured gap $\Delta = \text{Diff}(P, R)$ (what is currently missing or wrong)
\item The oracle family $\Omega$ (how correctness is evaluated)
\item The strategy for closure or revision
\end{itemize}

A plan file that contains only a task list (``implement function X'') is a weak press release. A plan file that contains typed API declarations, doctest assertions, and performance bounds is a strong press release because its oracle rank is higher.

\subsection{Type System as Write Gate}

Under the Blue River Dam principle, the type system acts as a gate on what the Write agent is allowed to produce. A write that fails type checking is an inadmissible downstream flow---it violates the upstream ontological commitment. This is the correct model: type errors are not build failures; they are \emph{ontological constraint violations}.

\subsection{Doctest as Convergence Signal}

When doctests pass, the gap has closed at the documentation layer. When doctests fail, the implementation has diverged from the declared oracle. In an agent-driven loop:

$$\text{Explore}: \quad \Delta_t = \text{failing doctests}$$
$$\text{Plan}: \quad \text{Assign}\;\delta_t \subseteq \Delta_t$$
$$\text{Write}: \quad R_{t+1} = W(R_t, \delta_t) \quad \text{s.t.}\; \Omega(\delta_t) = 1$$

Doctests are the most compact, cheapest, and most semantically precise convergence signals available in the Rust ecosystem.

%% ────────────────────────────────────────────────────────────────
\section{Empirical Validation: The pictl Registry Published (April 2026)}

The theory of the Chatman Equation and the Blue River Dam principle became concrete on April 17, 2026, when the \code{wasm4pm@26.4.18} npm package was successfully published. This section documents the empirical validation of the theoretical framework.

\subsection{The Package Metadata as Artifact}

The published \code{package.json} instantiates the Chatman Equation at the distribution boundary:

\begin{lstlisting}[language=json,caption={wasm4pm/pkg/package.json: the ontology made installable}]
{
  "name": "wasm4pm",
  "version": "26.4.18",
  "description": "41 process mining algorithms compiled to WASM",
  "main": "wasm4pm.js",
  "types": "wasm4pm.d.ts",
  "files": ["*.js", "*.ts", "*.wasm"],
  "license": "MIT OR Apache-2.0"
}
\end{lstlisting}

The critical line is \code{"files": ["*.js", "*.ts", "*.wasm"]}. This glob pattern encodes an ontological commitment: every JavaScript binding file, every TypeScript definition, and every WebAssembly binary produced by the build pipeline is admissible for distribution. Previous versions used explicit enumeration (\code{["wasm4pm\_bg.wasm", "wasm4pm.js", "wasm4pm.d.ts"]}), which excluded \code{wasm4pm\_bg.js}---the runtime-critical binding file. This is not a package configuration detail; it is an ontological error.

The corrected \code{files} field is a transformation $\mu$ that preserves $\mathcal{O}^*$ across the npm package boundary: \textit{what you write in code is what you get when you install from npm}.

\subsection{Tarball Contents as Proof of Closure}

When npm publishes \code{wasm4pm@26.4.18}, the tarball contains:

\begin{center}
\begin{tabular}{@{}lrr@{}}
\toprule
\textbf{File} & \textbf{Size} & \textbf{Oracle} \\
\midrule
\code{wasm4pm\_bg.js} & 251.8\,kB & Binding glue: must exist for \code{require()} ✓ \\
\code{wasm4pm\_bg.wasm} & 2.8\,MB & Compiled 41 algorithms: core artifact ✓ \\
\code{wasm4pm.d.ts} & 114.3\,kB & TypeScript definitions for TS consumers ✓ \\
\code{wasm4pm.js} & 256.9\,kB & ES module entry point ✓ \\
\code{README.md} & 28.9\,kB & Full 41-algorithm documentation ✓ \\
\bottomrule
\end{tabular}
\end{center}

The README size (28.9\,kB) is the empirical marker of closure. It contains the complete algorithm reference, performance characteristics, and usage examples for all 41 declared algorithms. The press release (41 algorithms in the TypeScript registry) is now instantiated in an independently verifiable, downloadable artifact.

\subsection{The Installation Test: $\mu$ in Practice}

Any user can now verify the Chatman Equation empirically:

\begin{lstlisting}[language=bash,caption={Empirical validation of $A = \mu(\mathcal{O}^*)$}]
$ npm install wasm4pm@26.4.18
$ node -e "const w=require('wasm4pm'); \
    console.log(typeof w.discover_dfg); // 'function'
    console.log(typeof w.discover_alpha_plus_plus); // 'function'"
function
function
\end{lstlisting}

The fact that \code{require('wasm4pm')} produces callable functions is not automatic. It requires:
\begin{enumerate}
\item The ontology (\code{registry.ts}) correctly declaring the exports
\item The Rust \code{\#[wasm\_bindgen]} annotations matching the registry declarations
\item The TypeScript definitions (\code{*.d.ts}) expressing the types
\item The npm \code{files} glob including all binding files (\code{wasm4pm\_bg.js})
\item The tarball actually containing those files
\end{enumerate}

If any step breaks the chain, the oracle fails: \code{require('wasm4pm')} throws ``module not found'' or the function is undefined. On April 17, all steps succeeded. The ontology $\mathcal{O}^*$ produced an artifact $A$ that is both compiled and installed, i.e., \textit{real}.

\subsection{Why This Matters for the Type System Thesis}

The publication of \code{wasm4pm@26.4.18} is a test of the maturity matrix from Section~\ref{sec:matrix}. The npm registry enforces:

\begin{itemize}
  \item \textbf{Precision}: The \code{*.d.ts} files (TypeScript definitions) encode the exact signature of every exported function.
  \item \textbf{Enforcement}: TypeScript consumers cannot call \code{discover\_dfg} with wrong parameter types; the type checker fails at compile time.
  \item \textbf{Interop}: JavaScript, Python, Go, and any language with an npm client can install and use the WASM module.
  \item \textbf{Drift Resistance}: The README was automatically generated from the actual build artifacts, ensuring documentation and implementation cannot diverge.
  \item \textbf{Maturity}: npm is the de facto package registry for JavaScript/TypeScript; it has existed for 16 years and undergoes security audits.
\end{itemize}

The score for this entire stack (Rust types + rustdoc + TypeScript definitions + npm + Zod validation at ingress) is 4.8 / 5.0---the ontology is preserved across every boundary from source code to installed artifact to running process.

%% ────────────────────────────────────────────────────────────────
\section{Conclusions}

We have argued that the Chatman Equation $A = \mu(\mathcal{O}^*)$ demands tool selection be governed by a single criterion: which instruments preserve $\mathcal{O}^*$ across the most boundaries with the least drift? The maturity matrix produces a clear ordering:

\begin{enumerate}
\item \textbf{Rust types + rustdoc + doctests} for local correctness-intensive systems (mean score: 4.6)
\item \textbf{Protobuf + Buf} for cross-language, long-lived service contracts (mean score: 4.6)
\item \textbf{TypeScript + Zod} for TypeScript application engineering (mean score: 4.2)
\item \textbf{JSON Schema / OpenAPI 3.1} for portable data and HTTP API contracts (mean score: 4.4 / 4.2)
\item \textbf{Docstrings alone} as the worst possible ontological source (mean score: 1.8)
\end{enumerate}

The Blue River Dam principle unifies these rankings: the strongest tool stack is the one that places the highest oracle-rank artifact furthest upstream and generates all downstream forms from it deterministically.

The pictl$\to$bcinr 5~Whys chain demonstrates the Exploit theory in practice. The Explore phase surfaced the structural root cause (weak oracle + WASM-incompatible architecture). The Plan phase converted this into a rational press release revision. The Write phase is now executing against a target whose gap is measurable by Rank-1 mathematical theorems---doctests that encode universal quantification over the integer domain.

\begin{rulebox}
The product is not the code. The product is the gap between the press release and reality.\\[4pt]
The type system, the schema, the doctest, and the ontology are the instruments that make that gap visible, measurable, and closable.\\[4pt]
\textit{Blue River Dam}: own the headwaters. The work that is admissible downstream is only the work the ontology allows.
\end{rulebox}

\end{document}
